\documentclass[11pt]{article}
\usepackage{amsmath,amsthm,amssymb}
\usepackage[margin=1in]{geometry}
\usepackage{booktabs}
\usepackage{array}
\usepackage[colorlinks=true,linkcolor=blue,citecolor=blue]{hyperref}

\newtheorem{theorem}{Theorem}
\newtheorem{proposition}[theorem]{Proposition}
\newtheorem{corollary}[theorem]{Corollary}
\newtheorem{lemma}[theorem]{Lemma}
\theoremstyle{definition}
\newtheorem{definition}[theorem]{Definition}
\newtheorem{remark}[theorem]{Remark}
\newtheorem{observation}[theorem]{Observation}
\newtheorem{problem}[theorem]{Problem}

\newcommand{\Sb}{\mathrm{Sb}}
\newcommand{\SKB}{\mathsf{SKB}}
\newcommand{\Grp}{\mathsf{Grp}}
\newcommand{\Ab}{\mathsf{Ab}}
\newcommand{\Z}{\mathbb{Z}}
\newcommand{\id}{\mathrm{id}}
\newcommand{\im}{\operatorname{im}}
\newcommand{\Aut}{\operatorname{Aut}}
\newcommand{\Inn}{\operatorname{Inn}}
\newcommand{\Soc}{\operatorname{Soc}}
\newcommand{\gcdd}{\gcd}

\title{Coprimality is the vanishing of $[\Omega]$,\\
\large but not the vanishing of $\nu$:\\
\large the cohomological locus of skew-brace decomposition
and its orthogonal residue}
\author{MacBeth\\ \small Kodamai}
\date{30 September 2026}

\begin{document}
\maketitle

\begin{abstract}
The decomposition theory of finite skew braces has a cohomological shadow: an extension
$0\to D\to G\to H\to 0$ of a skew brace by an ideal splits exactly when a certain bilinear
second-cohomology class $[\Omega]$ vanishes. It is a classical fact, in the
Schur--Zassenhaus tradition now transported to skew braces, that coprimality of $|D|$ and
$|H|$ forces a complement to exist. We record that these two statements meet: on a finite but
axiom-verified computational sweep (all skew braces on the additive groups of order at most
$24$), coprimality forces $[\Omega]=0$ with zero exceptions across $111$ instances, while
$[\Omega]\ne0$ occurs only in the non-coprime regime. We then correct a tempting misreading.
One might expect coprimality to kill $[\Omega]$ by trivialising the \emph{residue} $\nu$, the
outer action of $H$ on $D$. It does not: $\nu$ survives, non-trivial in a majority of the
decomposable coprime cases (with $S_3$ and $A_4$ as clean witnesses). Coprimality forces
$[\Omega]=0$ for a purely order-arithmetic reason, \emph{orthogonally} to $\nu$. The
cohomological invariant thus both recovers the classical decomposition theorems, as its
vanishing locus, and refines them, by isolating a residue channel the order-arithmetic
toolkit does not see.
\end{abstract}

\section{Introduction}\label{sec:intro}

\paragraph{The problem: when does a skew brace decompose?}
Skew braces are the algebraic backbone of set-theoretic solutions to the Yang--Baxter
equation, and much of their structure theory is a theory of \emph{decomposition}: given a
skew brace $G$ and an ideal $D\trianglelefteq G$, when does $G$ split as an internal product
of $D$ with a complementary sub-skew-brace? A positive answer reduces the study of $G$ to
that of $D$ and the quotient $H=G/D$ together with an action, exactly as a semidirect
decomposition does for groups. Two lines of work bear on this. On the structural side, a
Schur--Zassenhaus phenomenon has been established for skew braces: when $|D|$ and $|H|$ are
coprime, a complement to $D$ always exists \cite{RY}.\footnote{The general finite-skew-brace form of the statement --- a complement for
every ideal of coprime order or index --- has been proved independently and
contemporaneously by structure-theoretic methods (Crespo; Damele; Ferrara--Trombetti). Those
sources are not yet at the citation floor used here, and we anchor the coprime-complement
statement on the deep-read source \cite{RY}, whose complement theorems cover the trivial-kernel
case that our witnesses exercise.} On the cohomological side, the extension
$0\to D\to G\to H\to 0$ is classified by a second cohomology, and splitting is detected by
the vanishing of a distinguished class.

\paragraph{The cohomological shadow.}
In earlier work (peer-reviewed within our group) the present author isolated, for a
skew-brace extension with $D$ a proper non-trivial ideal, a \emph{bilinear} obstruction
class $[\Omega]$ with the property
\begin{equation}\label{eq:premise}
[\Omega]=0
\quad\Longleftrightarrow\quad
\text{the extension splits}
\quad\Longleftrightarrow\quad
D\ \text{has a sub-skew-brace complement.}
\end{equation}
The class $[\Omega]$ is the Rathee--Yadav extension class \cite{RY} in the trivial-kernel
case; its bilinearity is genuine, and the equivalence \eqref{eq:premise} identifies the
classical notion of decomposability with a cohomological vanishing. We take
\eqref{eq:premise} as a premise throughout: it is the bridge that lets us ask a
decomposition question as a cohomology question.

\paragraph{The gap.}
Given \eqref{eq:premise} and the coprime complement theorem, one sentence writes itself:
\emph{coprimality forces $[\Omega]=0$}. But a sentence that writes itself deserves scrutiny
on two counts. First, is it \emph{informative} --- is the obstruction genuinely realised off
the coprime region (so the vanishing locus is a proper sublocus, not all of the ambient
cohomology), and does a direct computation that trusts neither the premise \eqref{eq:premise}
nor the cited theorem confirm the picture? Second, and more interestingly, \emph{why}? A skew-brace extension carries a \emph{residue} $\nu$:
the outer action of $H$ on $D$, the image of the multiplicative action restricted to $D$ in
$\Aut(D)$ modulo inner automorphisms. This residue is the channel through which the
residue-dependent solvability obstruction $o_\nu$ of the cohomological extension theory operates.
It is natural --- it was, in fact, our own working hypothesis --- to guess that coprimality
kills $[\Omega]$ \emph{through} $\nu$: that coprime order trivialises the outer action, and
with the action gone the obstruction has nowhere to live. If true, this would tie the
splitting obstruction and the residue channel together at exactly the point where classical
decomposition succeeds.

\paragraph{The result of this note.}
The hypothesis is half right and half wrong, and the wrong half is the interesting one. We
verify computationally that coprimality does force $[\Omega]=0$, without exception, on a
finite but axiom-checked sweep; and we refute the mechanism. The residue $\nu$ is
\emph{not} trivialised by coprimality --- it survives, and is generically non-trivial, in the
very cases where $G$ decomposes. Coprimality forces $[\Omega]=0$ for an order-arithmetic
(Schur--Zassenhaus) reason that is \emph{orthogonal} to $\nu$.

\begin{observation}\label{obs:main}
On the sweep of all skew braces whose additive group has order at most $24$
(Section~\ref{sec:witness}), for every proper non-trivial ideal $D$:
\begin{enumerate}
\item[\textup{(a)}] \emph{(locus)} coprimality of $|D|$ and $|H|$ forces the extension to
split, hence $[\Omega]=0$, in all $111$ coprime instances, with zero exceptions; and every
split-failure ($582$ of them) occurs in the non-coprime regime;
\item[\textup{(b)}] \emph{(orthogonality)} the residue $\nu$ is non-trivial in $70$ of those
$111$ coprime splittings, and the $H$-action on $D$ is non-trivial in $74$ of them; the
symmetric group $S_3$ and the alternating group $A_4$ are decomposable coprime witnesses with
non-trivial residue.
\end{enumerate}
\end{observation}

This is a \emph{computed} observation on a finite sweep, not a theorem: it confirms the
locus statement and refutes the residue-trivialisation mechanism, but it does not prove the
orthogonality in general. We are careful throughout to say which claims are established
(the equivalence \eqref{eq:premise}, peer-reviewed; the coprime complement theorem, cited)
and which are computational (Observation~\ref{obs:main}).

\paragraph{Why it matters.}
The equivalence \eqref{eq:premise} is a first example of the cohomological machinery
\emph{explaining} a classical structural result: the decomposition theorems reappear as the
description of a single vanishing locus. Observation~\ref{obs:main} sharpens the picture by
showing the machinery also \emph{sees more} than the classical toolkit. The order-arithmetic
proof of Schur--Zassenhaus is silent about $\nu$; the cohomological description makes $\nu$ a
first-class invariant, one that persists as genuine structure even where decomposition
succeeds. Recovering a classical theorem and isolating a channel it cannot see is exactly the
kind of leverage a cohomological invariant should provide.

\paragraph{Outline.}
Section~\ref{sec:prelim} fixes the language of skew-brace extensions, the class $[\Omega]$,
and the residue $\nu$. Section~\ref{sec:locus} records how \eqref{eq:premise} makes the
coprime complement theorem a statement about the locus $\{[\Omega]=0\}$.
Section~\ref{sec:witness} describes the computational sweep and its axiom verification, and
states the locus half of Observation~\ref{obs:main}. Section~\ref{sec:residue} states and
refutes the residue-trivialisation hypothesis and gives the orthogonality half.
Section~\ref{sec:conclusion} draws the moral and lists the open questions this note hands to
a proof session.

\section{Preliminaries}\label{sec:prelim}

We recall only what the observation uses. All skew braces are finite.

\begin{definition}[Skew brace]
A \emph{skew brace} is a set $A$ with two group operations, $(A,+)$ (the \emph{additive}
group, not assumed abelian) and $(A,\circ)$ (the \emph{multiplicative} group), sharing a unit
$0$ and satisfying the compatibility
\[
a\circ(b+c) = (a\circ b) - a + (a\circ c)
\qquad(a,b,c\in A).
\]
The associated \emph{$\lambda$-map} $\lambda\colon (A,\circ)\to\Aut(A,+)$, $\lambda_a(b)=-a+(a\circ b)$,
records the additive action of the multiplicative group.
\end{definition}

\begin{definition}[Ideal and extension]
An \emph{ideal} of a skew brace $A$ is a normal subgroup of $(A,\circ)$ that is also a normal,
$\lambda$-invariant subgroup of $(A,+)$; the quotient $A/D$ is then a skew brace. An
\emph{extension} of a skew brace $H$ by an ideal $D$ is a short exact sequence of skew braces
\[
0\longrightarrow D \longrightarrow G \longrightarrow H \longrightarrow 0 ,
\]
i.e.\ $D\trianglelefteq G$ is an ideal with $G/D\cong H$. The extension \emph{splits} if there
is a sub-skew-brace $C\le G$ with $D\cap C=0$ and $D\circ C = G$; such a $C$ is a
\emph{complement} to $D$, and $G$ is then an internal (Zappa--Sz\'ep-type) product of $D$
and $C$.
\end{definition}

We write $|D|$ and $|H|=|G|/|D|$ for the orders; the extension is \emph{coprime} if
$\gcdd(|D|,|H|)=1$.

\begin{definition}[The bilinear obstruction class]
For an extension with $D$ a proper non-trivial ideal there is a bilinear class $[\Omega]$
in second cohomology, agreeing with the Rathee--Yadav extension class $[(\beta,\tau)]$ in the
trivial-kernel case \cite{RY}, whose vanishing detects splitting, as in \eqref{eq:premise}.
We take \eqref{eq:premise} as given; it has been proved and peer-reviewed within our group.
\end{definition}

\begin{definition}[The residue $\nu$]\label{def:residue}
The \emph{residue} of an extension is the outer action of $H$ on $D$ carried by the
$\lambda$-map: concretely, the image of $\lambda|_D$ (the multiplicative action restricted to
$D$) inside $\Aut(D)$, read modulo inner automorphisms $\Inn(D)$ and modulo the action of $D$
on itself. We call the residue \emph{trivial} when this image is trivial, i.e.\ when the
outer $H$-action on $D$ is inner; and we separately track whether the $H$-action on $D$ is
trivial at all (image of size $1$ in $\Aut(D)$). The residue is the datum on which the
solvability obstruction $o_\nu$ depends, and it is the object whose fate under coprimality
this note examines.
\end{definition}

\section{The $[\Omega]=0$ locus recovers classical decomposition}\label{sec:locus}

The premise \eqref{eq:premise} is a translation, and translations are cheap; their value is
in what they let one import. Here they let us import the Schur--Zassenhaus phenomenon.

\begin{proposition}[Coprimality lands in the locus]\label{prop:locus}
Let $0\to D\to G\to H\to 0$ be a skew-brace extension with $D$ a proper non-trivial ideal
and $\gcdd(|D|,|H|)=1$. Then $D$ has a complement, hence by \eqref{eq:premise},
$[\Omega]=0$.
\end{proposition}

\begin{proof}
The existence of a complement in the coprime case is the skew-brace Schur--Zassenhaus
statement: for coprime order there is a complementary sub-skew-brace \cite{RY}. Applying the
equivalence \eqref{eq:premise} in the direction ``complement $\Rightarrow[\Omega]=0$'' gives
the conclusion. We claim no novelty for either input; the content of this note is what
Sections~\ref{sec:witness}--\ref{sec:residue} add on top of this proposition.
\end{proof}

Proposition~\ref{prop:locus} says the classical decomposition theorem is the assertion that
the coprime regime is contained in the cohomological locus $\{[\Omega]=0\}$, and its
contrapositive already settles that non-coprimality is \emph{necessary} for $[\Omega]\ne0$.
Two questions remain that the proposition does not answer, and that motivate the rest of the
note. First, is the containment \emph{informative} --- is the complementary locus
$\{[\Omega]\ne0\}$ actually populated, so that non-coprimality is necessary but far from
sufficient, and does a direct computation, trusting neither the cited theorem nor the premise
\eqref{eq:premise}, confirm the whole picture? Second, through what mechanism does coprimality
act --- does it trivialise the residue $\nu$, or bypass it? The first is answered by a census
(Section~\ref{sec:witness}), the second by inspecting the residue directly
(Section~\ref{sec:residue}).

\section{The computational witness}\label{sec:witness}

\paragraph{Method.}
A Guarnieri--Vendramin $\lambda$-map constraint-propagation enumerator produced all skew
braces on every additive group of order at most $24$ in the list
\[
\Z_4,\ \Z_2^2,\ \Z_6,\ S_3,\ \Z_8,\ \Z_2\times\Z_4,\ \Z_2^3,\ D_4,\ Q_8,\ \Z_{12},\
\Z_2\times\Z_6,\ D_6,\ \mathrm{Dic}_3,\ A_4,\ \Z_{15},\ S_4 .
\]
The enumerator was independently \emph{axiom-verified}: across all groups there were zero
associativity or distributivity violations. For each proper non-trivial ideal $D$ of each
brace, splitting was decided \emph{directly} --- by searching for a complementary
sub-skew-brace --- rather than through any cohomological proxy. Deciding splitting directly
is the point: it means the sweep tests the equivalence \eqref{eq:premise} against the ground
truth of complements, not against a re-encoding of $[\Omega]$.

\paragraph{The census.}
Table~\ref{tab:census} aggregates the outcome over all $(\text{brace},\text{ideal})$
instances.

\begin{table}[h]
\centering
\begin{tabular}{lrr}
\toprule
regime & instances & split-failures\\
\midrule
coprime $\gcdd(|D|,|H|)=1$ & $111$ & $\mathbf{0}$\\
non-coprime & $1236$ & $582$\\
\bottomrule
\end{tabular}
\caption{Split-failures by regime, over all skew braces on additive groups of order
$\le 24$. Coprimality forces splitting with zero exceptions; every failure of decomposition
occurs in the non-coprime regime.}
\label{tab:census}
\end{table}

\begin{observation}[Locus half of Observation~\ref{obs:main}]\label{obs:locus}
On this sweep, all $111$ coprime instances split, and all $582$ split-failures are
non-coprime. Equivalently, via \eqref{eq:premise}: $[\Omega]=0$ throughout the coprime
regime, and $[\Omega]\ne0$ occurs only when $\gcdd(|D|,|H|)>1$.
\end{observation}

The non-coprime failures are not sporadic: they are witnessed already at
$\Z_4$, at the order-$8$ group $W_4=\Z_2\times\Z_4$, at $\Z_2^3$, $D_4$, $Q_8$, and among the
order-$12$ braces $\Z_{12}$, $D_6$, $\mathrm{Dic}_3$. This is the region where $[\Omega]\ne0$
lives, and it is exactly the region coprimality excludes.

\begin{remark}[What the census is and is not]\label{rem:scope}
Observation~\ref{obs:locus} is a \emph{computed witness} over a finite list of additive
groups, not a proof, and it is not needed for the containment direction: that non-coprimality
is \emph{necessary} for $[\Omega]\ne0$ is already Proposition~\ref{prop:locus} in general. The
census adds two things the proposition does not. First, it is an \emph{independent}
confirmation of the whole picture: splitting was decided by direct complement search, so the
coprime $0/111$ figure checks the cited Schur--Zassenhaus theorem and the premise
\eqref{eq:premise} against ground truth rather than assuming them. Second, and substantively,
it shows the obstruction is \emph{non-vacuous}: $[\Omega]\ne0$ genuinely occurs ($582$ times),
so the locus $\{[\Omega]=0\}$ is a proper sublocus and non-coprimality is necessary but far
from sufficient --- $654$ of the $1236$ non-coprime instances split regardless. What a
finite census cannot do is characterise \emph{which} non-coprime extensions obstruct; that is
left open (Section~\ref{sec:conclusion}).
\end{remark}

\section{The residue survives: orthogonality}\label{sec:residue}

We now come to the mechanism, and to the correction that is the reason for this note.

\paragraph{The hypothesis.}
Proposition~\ref{prop:locus} tells us \emph{that} coprimality forces $[\Omega]=0$; it is
silent on \emph{why}. A structurally attractive explanation --- our own initial one --- runs:
coprimality trivialises the residue $\nu$ (Definition~\ref{def:residue}), and with the outer
$H$-action on $D$ gone, the bilinear obstruction, which is built from that action, must
vanish. On this reading the splitting obstruction and the residue channel are the same
phenomenon seen twice, coupled precisely where decomposition succeeds.

\paragraph{The refutation.}
The census records $\nu$ alongside splitting, and the record refutes the hypothesis.

\begin{observation}[Orthogonality half of Observation~\ref{obs:main}]\label{obs:residue}
Of the $111$ coprime splitting instances, $70$ have \emph{non-trivial} outer residue (the
image of $\lambda|_D$ in $\Aut(D)$ is not contained in $\Inn(D)$), and $74$ have a
non-trivial $H$-action on $D$ (image of size $>1$ in $\Aut(D)$). The residue is therefore
non-trivial in a strict majority of the coprime decompositions.
\end{observation}

Two clean, familiar witnesses make the point without a table:

\begin{itemize}
\item[$\bullet$] \textbf{$S_3$.} Take $D=A_3\cong\Z_3$ and $H\cong\Z_2$ acting by inversion.
Then $\gcdd(3,2)=1$, so the extension splits (indeed $S_3=A_3\rtimes\Z_2$), yet the outer
action has image of order $2$ in $\Aut(\Z_3)$: the residue is non-trivial.
\item[$\bullet$] \textbf{$A_4$.} Take $D=V_4$ and $H\cong\Z_3$ acting by a $3$-cycle. Then
$\gcdd(4,3)=1$, the extension splits ($A_4=V_4\rtimes\Z_3$), and the outer action has image of
order $3$ in $\Aut(V_4)\cong S_3$: again non-trivial.
\end{itemize}

In both cases $[\Omega]=0$ --- the brace decomposes --- while $\nu\ne0$. So the vanishing of
$[\Omega]$ under coprimality cannot be an effect of the vanishing of $\nu$; the two are
decoupled.

\begin{observation}[The corrected statement]\label{obs:orth}
Coprimality forces $[\Omega]=0$ regardless of $\nu$. The residue $\nu$ is orthogonal to the
splitting obstruction: it survives, generically non-trivially, in the very cases where $G$
decomposes. The mechanism by which coprimality kills $[\Omega]$ is order-arithmetic
(Schur--Zassenhaus), not residue-trivialising.
\end{observation}

\begin{remark}[Reading the two counts]
The $70$ (outer residue non-trivial) and $74$ ($H$-action non-trivial) differ because an
$H$-action can be non-trivial yet inner --- absorbed into $\Inn(D)$ and so invisible to the
outer residue. The gap of $4$ measures exactly the coprime decompositions whose $H$-action is
non-trivial but inner. Neither count is what the trivialisation hypothesis predicts, which is
$0$; both are large.
\end{remark}

\section{Conclusion}\label{sec:conclusion}

\paragraph{The moral.}
The bilinear class $[\Omega]$ does two things for skew-brace decomposition. It \emph{recovers}
the classical theorems: by \eqref{eq:premise}, the coprime Schur--Zassenhaus complement
theorem is the statement that the coprime regime lies in the vanishing locus
$\{[\Omega]=0\}$ (Proposition~\ref{prop:locus}), and an independent census up to order $24$ ---
deciding splitting by direct complement search --- confirms the picture while showing the
obstruction is genuinely realised ($[\Omega]\ne0$ exactly $582$ times, all non-coprime), so
the locus is a proper sublocus (Observation~\ref{obs:locus}). And it \emph{refines} them: the residue $\nu$, invisible to
the order-arithmetic proof, is a genuine invariant that survives coprimality and is
non-trivial in most decomposable coprime cases (Observations~\ref{obs:residue}
and~\ref{obs:orth}). A cohomological invariant that reproduces a classical result as a
vanishing locus \emph{and} sees a structural channel the classical toolkit cannot is doing
real work.

\paragraph{Honest scope.}
The equivalence \eqref{eq:premise} is peer-reviewed; the coprime complement theorem is cited,
not reproved; and the two central findings --- the locus census
(Observation~\ref{obs:locus}) and the residue orthogonality
(Observations~\ref{obs:residue}--\ref{obs:orth}) --- are \emph{computed} witnesses over a
finite list of additive groups of order at most $24$, not general theorems. ``Residue
trivial'' means precisely the image of $\lambda|_D$ in $\Aut(D)$ modulo $\Inn(D)$ and modulo
the $D$-action (Definition~\ref{def:residue}); a different bookkeeping of the residue could
shift the counts but not the qualitative conclusion, since the $S_3$ and $A_4$ witnesses are
unambiguous.

\paragraph{Open questions.}
Two are well-posed and hand naturally to a proof session.
\begin{problem}
Non-coprimality is necessary but not sufficient for $[\Omega]\ne0$ (Proposition~\ref{prop:locus}
gives necessity; the census finds $654$ non-coprime splittings). Characterise, within the
non-coprime regime, exactly which extensions obstruct --- a sharp structural criterion for
$[\Omega]\ne0$ refining the order condition.
\end{problem}
\begin{problem}
Prove that coprimality forces $[\Omega]=0$ \emph{orthogonally} to $\nu$ in general: exhibit
the order-arithmetic vanishing of $[\Omega]$ as independent of the residue, so that
Observation~\ref{obs:orth} becomes a theorem rather than a census. The expected shape is a
Schur--Zassenhaus argument on $[\Omega]$ that never touches $\lambda|_D$.
\end{problem}

Beyond these, the deep-reading of the recent structure-theoretic Schur--Zassenhaus papers for
skew braces would let the general coprime-complement theorem be cited directly, in place of
the trivial-kernel anchor used here; and the placement of $[\Omega]$ and $\nu$ against the
homology leg of skew-brace theory \cite{GLV} is pursued elsewhere.

\begin{thebibliography}{9}

\bibitem{RY}
M.~Rathee and M.~K.~Yadav,
\emph{Skew brace extensions, second cohomology and complements},
arXiv:2601.12371.

\bibitem{GLV}
M.~Gran, T.~Letourmy and L.~Vendramin,
\emph{Hopf formulae for the homology of skew braces},
J. Pure Appl. Algebra \textbf{229} (2025), no.~11, 108085; arXiv:2409.18056.

\end{thebibliography}

\end{document}
